% ── CODIFICAÇÃO E IDIOMA ──────────────────────────────────────────────────────

% Codificação UTF-8 para caracteres portugueses (ã, ç, ê, etc.)
\usepackage[utf8]{inputenc}

% Codificação T1 para hifenização correta de palavras acentuadas
\usepackage[T1]{fontenc}

% Regras de hifenização e formatação de datas em português brasileiro
\usepackage[brazilian]{babel}

% ── TIPOGRAFIA ────────────────────────────────────────────────────────────────

% Times New Roman para texto corrido e modo matemático — padrão UFN
\usepackage{mathptmx}

% ── GEOMETRIA (MARGENS UFN) ───────────────────────────────────────────────────

% Sobrescreve os padrões do IEEEtran para as margens de submissão da UFN
\usepackage[
  a4paper,
  top=2.5cm,
  bottom=2.5cm,
  left=1.8cm,
  right=1.8cm
]{geometry}

% ── FIGURAS ───────────────────────────────────────────────────────────────────

% \includegraphics — figuras simples da pasta /images
\usepackage{graphicx}

% \begin{subfigure} — figuras lado a lado (ex.: vistas do braço robótico)
\usepackage{subcaption}

% ── MATEMÁTICA ────────────────────────────────────────────────────────────────

% Ambientes align, equation, matrix — para fórmulas de IK/IA
\usepackage{amsmath}

% \mathbb{R}, \forall, \exists — para notação formal
\usepackage{amssymb}

% ── CITAÇÕES ──────────────────────────────────────────────────────────────────

% Compressão de citações numéricas estilo IEEE: [1]–[4] em vez de [1][2][3][4]
\usepackage{cite}

% ── LISTAGENS DE CÓDIGO ───────────────────────────────────────────────────────

% Pacotes para exibição de código-fonte com numeração de linhas
\usepackage{listings}
\usepackage{xcolor}

% Configuração de listagens para código monoespaçado padrão UFN
\lstset{
  basicstyle=\ttfamily\footnotesize,
  numbers=left,
  numberstyle=\tiny\color{gray},
  stepnumber=1,
  numbersep=5pt,
  backgroundcolor=\color{gray!10},
  showspaces=false,
  showstringspaces=false,
  showtabs=false,
  frame=single,
  rulecolor=\color{black!30},
  tabsize=2,
  captionpos=b,
  breaklines=true,
  breakatwhitespace=false,
  keywordstyle=\color{blue!70}\bfseries,
  commentstyle=\color{green!50!black}\itshape,
  stringstyle=\color{red!60},
  language=Python,   % padrão; substituir por language=Java, C++, etc. por listagem
  literate=
    {á}{{\'a}}1 {é}{{\'e}}1 {í}{{\'i}}1 {ó}{{\'o}}1 {ú}{{\'u}}1
    {Á}{{\'A}}1 {É}{{\'E}}1 {Í}{{\'I}}1 {Ó}{{\'O}}1 {Ú}{{\'U}}1
    {ã}{{\~a}}1 {õ}{{\~o}}1 {ç}{{\c{c}}}1 {ñ}{{\~n}}1
}
% O bloco literate trata caracteres portugueses dentro de blocos de código

% ── TABELAS ───────────────────────────────────────────────────────────────────

% Colunas auto-expansíveis — necessário para a tabela de cronograma do TCC 2
\usepackage{tabularx}

% \toprule, \midrule, \bottomrule — réguas profissionais (sem \hline)
\usepackage{booktabs}

% ── METADADOS PDF E HIPERLINKS ────────────────────────────────────────────────

% Deve ser o ÚLTIMO pacote carregado — gerencia metadados e referências clicáveis
\usepackage[
  pdftitle={TCC 1 --- [TÍTULO DO TRABALHO]},
  pdfauthor={[NOME DO ALUNO]},
  pdfsubject={Trabalho de Conclusão de Curso --- UFN},
  pdfkeywords={[palavra1], [palavra2], [palavra3]},
  colorlinks=true,
  linkcolor=black,
  citecolor=black,
  urlcolor=blue!70!black,
  bookmarks=true
]{hyperref}

% ── COMANDOS PERSONALIZADOS ABNT ──────────────────────────────────────────────

% ABNT NBR 10520 — Citação longa (> 3 linhas):
% Recuo de 4cm à esquerda, fonte \small, espaçamento simples, sem aspas
\newenvironment{citacaolonga}{%
  \begin{list}{}{%
    \setlength{\leftmargin}{4cm}%
    \setlength{\rightmargin}{0cm}%
    \setlength{\topsep}{0pt}%
    \setlength{\partopsep}{0pt}%
  }%
  \item\small%
}{%
  \end{list}%
}
% Uso: \begin{citacaolonga} texto da citação... \end{citacaolonga}

% Macro de conveniência para o nome da instituição
\newcommand{\UFN}{Universidade Franciscana (UFN)}
